\Article{22}
Da die in dieser Abhandlung dargestellte Methode besonders häufig und bequem bei Rechnungen der höheren Geodäsie angewandt wird, hoffen wir, dass die Leser eine Erläuterung der Regeln anhand einiger Beispiele aus diesem Gebiet begrüßen werden.

Die Bedingungsgleichungen zwischen den Winkeln eines Dreieckssystems sind hauptsächlich aus drei Quellen zu gewinnen.

\par\noindent\textbf{I.} Die Summe der Horizontalwinkel, die um denselben Scheitel einen vollen Horizontumlauf bilden, muss vier rechten Winkeln gleich sein.

\par\noindent\textbf{II.} Die Summe der drei Winkel in jedem Dreieck ist einer gegebenen Größe gleich, denn wenn das Dreieck auf einer gekrümmten Fläche liegt, lässt sich der Überschuss dieser Summe über zwei rechte Winkel so genau berechnen, dass er als absolut genau gelten kann.

\par\noindent\textbf{III.} Die dritte Quelle ist das Verhältnis der Seiten in Dreiecken, die eine geschlossene Kette bilden. Ist eine Reihe von Dreiecken so verbunden, dass das zweite Dreieck eine Seite \(a\) mit dem ersten und eine andere, \(b\), mit dem dritten gemeinsam hat; ebenso das vierte Dreieck eine Seite \(c\) mit dem dritten und eine Seite \(d\) mit dem fünften, und so fort bis zum letzten Dreieck, dessen gemeinsame Seite mit dem vorhergehenden Dreieck \(k\) und dessen weitere gemeinsame Seite mit dem ersten Dreieck wiederum \(l\) ist, so sind die Werte der Verhältnisse \(\frac al,\allowbreak\frac ba,\allowbreak\frac cb,\allowbreak\frac dc\ .\ .\ .\ .\ \frac lk\) nach bekannten Methoden jeweils aus den beiden Winkeln aufeinanderfolgender Dreiecke bekannt, die den gemeinsamen Seiten gegenüberliegen. Da das Produkt dieser Brüche \(=1\) sein muss, ergibt sich eine Bedingungsgleichung zwischen den Sinus dieser Winkel (wobei diese Winkel, falls die Dreiecke auf einer gekrümmten Fläche liegen, jeweils um den dritten Teil des sphärischen oder sphäroidischen Exzesses zu vermindern sind).

In komplizierteren Dreiecksnetzen kommt es jedoch häufig vor, dass mehr Bedingungsgleichungen zweiter und dritter Art auftreten, als beibehalten werden dürfen, weil ein Teil davon bereits in den übrigen enthalten ist. Seltener ist der Fall, dass den Bedingungsgleichungen zweiter Art ähnliche Gleichungen für Figuren mit mehr Seiten hinzugefügt werden müssen, nämlich nur dann; wenn \SourceLeaf{087}
Vielecke gebildet werden, die durch Messungen nicht in Dreiecke zerlegt sind. Auf diese Dinge, die dem gegenwärtigen Gegenstand zu fern liegen, werden wir bei anderer Gelegenheit ausführlicher eingehen. Einen Hinweis dürfen wir jedoch nicht übergehen: Soll unsere Theorie rein und streng angewandt werden, setzt sie voraus, dass die durch \(v,v',v''\) etc. bezeichneten Größen tatsächlich entweder unmittelbar beobachtet oder aus Beobachtungen so hergeleitet sind, dass sie voneinander unabhängig bleiben, oder wenigstens als unabhängig angesehen werden können. In der üblichen Praxis beobachtet man die Winkel der Dreiecke selbst; sie können daher als \(v,v',v''\) etc. angesetzt werden. Wir müssen jedoch beachten, dass bei einem System, das zusätzlich Dreiecke enthält, deren Winkel nicht unmittelbar beobachtet wurden, sondern als Summen oder Differenzen tatsächlich beobachteter Winkel entstehen, diese nicht zu den beobachteten Größen gezählt werden dürfen; vielmehr muss ihre zusammengesetzte Form in den Rechnungen beibehalten werden. Anders verhält es sich bei einer Beobachtungsweise ähnlich derjenigen des berühmten \textsc{Struve} (Astronomische Nachrichten II, S. 431), bei der die Richtungen der einzelnen Seiten, die von demselben Scheitel ausgehen, durch Vergleich mit ein und derselben willkürlichen Richtung bestimmt werden. Dann sind eben diese Winkel als \(v,v',v''\) etc. anzusetzen; dadurch treten alle Dreieckswinkel in Differenzenform auf, und die Bedingungsgleichungen erster Art, denen von Natur aus von selbst genügt wird, entfallen als überflüssig. Die Beobachtungsweise, die ich selbst bei der in den vorangegangenen Jahren abgeschlossenen Triangulation befolgt habe, unterscheidet sich zwar sowohl von der ersten als auch von der zweiten; in ihrer Wirkung lässt sie sich jedoch der zweiten gleichsetzen. Daher müssen an jeder Station die von dort ausgehenden Seitenrichtungen, die von einer gleichsam willkürlich gewählten Anfangsrichtung aus gezählt werden, als Größen \(v,v',v''\) etc. angesetzt werden. Wir werden nun zwei Beispiele durcharbeiten: eines gehört zur ersten, das andere zur zweiten Beobachtungsweise.
